\section{Лекция 14}
\begin{theorem}[Сферические координаты в многомерном пространстве] $\\$
    Определим $\phi$ по правилу
    \[\begin{cases}
        x_1=r\cdot \cos{\alpha_1}\cdot \dots \cdot \cos{\alpha_{n-1}}\\
        x_2=r\cdot \sin{\alpha_1}\cdot \cos{\alpha_2}\cdot \dots \cdot \cos{\alpha_{n-1}}\\
        x_3=r\cdot \sin{\alpha_2}\cdot \cos{\alpha_3}\cdot \dots \cdot \cos{\alpha_{n-1}}\\
        \vdots\\
        x_{n-1}=r\cdot \sin{\alpha_{n-2}}\cdot \cos{\alpha_{n-1}}\\
        x_n=r\cdot \sin{\alpha_{n-1}}
    \end{cases} \eqno{(*)}\]
    Тогда $\phi\in \mathcal{C}^1$, причем $\phi$ является биекцией между множеством
    \[\Omega=\{r>0,\ \alpha_1\in (0,2\pi),\ \alpha_j\in \left(-\frac{\pi}{2},\frac{\pi}{2}\right),\ j=2,\dots,n-1\}\]
    и множеством
    \[G=\R^n\setminus \{x_2=0, x_1\geq 0\}\]
    прчием 
    \[\det{J}_{\phi}=r^{n-1}\cos{\alpha_2}\cdot \cos^2{\alpha_3}\dots\cos^{n-2}{\alpha_{n-1}}\]
\end{theorem}
\begin{proof}
    Очевидно, что $\phi\in \mathcal{C}^1(\Omega)$. Заметим, что на $\Omega$ выполнено $\cos{\alpha_j}>0$ для $j=2,\dots,n-1$ и 
    \[x_2=0 \Leftrightarrow \sin{\alpha_1}=0 \Leftrightarrow \alpha_1=\pi \Rightarrow \cos{\alpha_1}=-1,\ x_1<0\]
    Таким образом $\phi(\Omega)\subset G$. Рассмотрим произвольную точку $(x_1,\dots,x_n)\in G$ Возведем уравнения $(*)$ в квадрат и последовательно сложим:
    \[\begin{cases}
    x_1^2=r^2\cdot \cos^2{\alpha_1}\cdot \cos^2{\alpha_2}\dots\cos^2{\alpha_{n-1}}\\
    x_1^2+x_2^2=r^2\cdot \cos^2{\alpha_2}\dots\cos^2{\alpha_{n-1}}\\
    \vdots\\
    x_1^2+\dots+x_{n-1}^2=r^2\cdot \cos^2{\alpha_{n-1}}\\
    x_1^2+\dots+x_n^2=r^2
    \end{cases}\]
    \[r=\sqrt{x_1^2+\dots+x_n^2}>0\]
    Для $j=2,\dots,n-1$ определим косинус:
    \[\cos{\alpha_j}=\frac{\sqrt{x_1^2+\dots+x_j^2}}{\sqrt{x_1^2+\dots+x_{j+1}^2}}>0\]
    а для $\alpha_1$:
    \[\cos{\alpha_1}=\frac{x_1}{\sqrt{x_1^2+x_2^2}},\ \sin{\alpha_1}=\frac{x_2}{\sqrt{x_1^2+x_2^2}}\]
    Поскольку $\alpha\in(0,2\pi)$, то $\alpha_1$ определен однозначно. Теперь для $j=2,\dots,n-1$ определим синус:
    \[\sin{\alpha_j}=\frac{x_{j+1}}{\sqrt{x_1^2+\dots+x_{j+1}^2}}\]
    Так как $\alpha_j\in \left(-\frac{\pi}{2}, \frac{\pi}{2}\right)$, то однозначно определен
    \[\alpha_j=\arcsin\left(\frac{x_{j+1}}{\sqrt{x_1^2+\dots+x_{j+1}^2}}\right)\]
    Остается проверить, что построенные $(r,\alpha_1,\dots,\alpha_{n-1})$ задают исходную точку:
    \begin{multline*}
        r\cdot\cos{\alpha_1}\cdot \dots\cdot \cos{\alpha_{n-1}}=\\
        =r\cdot \frac{x_1}{\sqrt{x_1^2+x_2^2}}\cdot \frac{\sqrt{x_1^2+x_2^2}}{\sqrt{x_1^2+x_2^2+x_3^2}}\cdot \dots \cdot \frac{\sqrt{x_1^2+\dots+x_{n-1}^2}}{r}=x_1
    \end{multline*}
    \begin{multline*}
        r\cdot\sin{\alpha_1}\cdot \cos{\alpha_2}\cdot \dots\cdot \cos{\alpha_{n-1}}=\\
        =r\cdot\frac{x_2}{\sqrt{x_1^2+x_2^2}}\cdot \frac{\sqrt{x_1^2+x_2^2}}{\sqrt{x_1^2+x_2^2+x_3^2}}\cdot \dots \cdot \frac{\sqrt{x_1^2+\dots+x_{n-1}^2}}{r}=x_2
    \end{multline*}
    Теперь для $j=2,\dots,n-1$:
    \begin{multline*}
        r\cdot \sin{\alpha_j}\cdot \cos{\alpha_{j+1}\cdot \dots \cdot \cos{\alpha_{n-1}}}=\\
        =r\cdot \frac{x_{j+1}}{\sqrt{x_1^2+\dots+x_{j+1}^2}}\cdot \frac{\sqrt{x_1^2+\dots+x_{j+1}^2}}{\sqrt{x_1^2+\dots+x_{j+2}^2}}\cdot \dots \cdot \frac{\sqrt{x_1^2+\dots+x_{n-1}^2}}{r}=x_{j+1}
    \end{multline*}
    Таким образом для любой точки $(x_1,\dots,x_n)\in G$ существует единственный набор $(r,\alpha_1,\dots,\alpha_{n-1})\in \Omega$, что $\phi(r,\alpha_1,\dots,\alpha_{n-1})=(x_1,\dots,x_n)$.
\end{proof}
%\noindent Пусть $\phi:\Omega\to G,\ \phi\in \mathcal{C}^1$ --- биекция и $|\det{J_{\phi}}|\neq 0$ 
%\[\idotsint\limits_{K}f(x_1,\dots,x_k)\ dx_1\dots dx_k=\idotsint\limits_{\Phi}f(\phi(y_1,\dots,y_k))\cdot |\det{J_{\phi}}|\ dy_1 \dots dy_k\]
\begin{theorem} [Теорема о замене переменных в несобственном кратном интеграле] $\\$
    Пусть $\phi: \Omega\to G,\ \phi\in \mathcal{C}^1(\Omega)$ --- биекция и $\det{J_{\phi}}$ не обращается в ноль на $\Omega$. Тогда, при условии что оба приведенных ниже несобственных интеграла существуют, то
    \[\idotsint\limits_{G}f(x_1,\dots,x_k)\ dx_1\dots dx_k=\idotsint\limits_{\Omega}f(\phi(y_1,\dots,y_k))\cdot |\det{J_{\phi}}|\ dy_1 \dots dy_k\]
\end{theorem}
\begin{proof}
    Пусть $\{\Phi_n\}$ --- исчерпание $\Omega$, тогда $\{K_n\}=\{\phi(\Phi_n)\}$ --- исчерпание $G$. По теореме \ref{zamena} о замене переменной в кратном интеграле
    \[\idotsint\limits_{K_n}f(x_1,\dots,x_k)\ dx_1\dots dx_k=\idotsint\limits_{\Phi_n}f(\phi(y_1,\dots,y_k))\cdot |\det{J_{\phi}}|\ dy_1 \dots dy_k\]
    поскольку оба интеграла сущетсвуют, то, устремив $n\to \infty$, получим утверждение теоремы.
\end{proof}
\begin{reminder}
    Последовательность $\{a_n\}$ называется фундаментальной, если 
    \[\forall \epsilon>0\ \exists\ N\in \N,\ \forall m,n>N: |a_n-a_m|<\epsilon\]
\end{reminder}
\begin{theorem}[Критерий Коши существования несобственного кратного интграла] $\\$
    Пусть $\Omega$ --- открытое множество, $f\in \mathcal{R}_{\text{loc}}(\Omega)$. Тогда $f\in \widetilde{\mathcal{R}}(\Omega)$ тогда и только тогда, когда $\forall \epsilon>0\ \exists\ K_{\epsilon}\subset \Omega$ --- измеримый компакт такой, что для любого измеримого компакта $K\subset\Omega\setminus K_{\epsilon}$:
    \[\idotsint\limits_{K}|f|\ dx_1 \dots dx_k<\epsilon\]
\end{theorem}
\begin{proof}\tab
    \begin{itemize}
        \item[$(\Rightarrow)$:]
        Пусть $f\in \widetilde{\mathcal{R}}(\Omega) \Rightarrow |f|\in \widetilde{\mathcal{R}}(\Omega)$. Пусть $\{K_n\}$ --- какое-либо исчерпание $\Omega$. Определим последовательность
        \[a_n=\idotsint\limits_{K_n}|f|\ dx_1 \dots dx_k\]
        она имеет предел, а значит она фундаментальна, то есть \\
        $\forall \epsilon>0\ \exists\ N\in \N,\ \forall m,n>N$:
        \[\idotsint\limits_{K_m}|f|\ dx_1 \dots dx_k-\idotsint\limits_{K_n}|f|\ dx_1 \dots dx_k<\epsilon\]
        Для каждого $\epsilon>0$ положим $K_{\epsilon}=K_{N+1}$. Пусть $K\subset \Omega\setminus K_{N+1}$. Поскольку $\{K_n\}$ --- исчерпание $\Omega$, то $\exists\ m \in \N: K\subset K_m$. Значит $K\subset K_m\setminus K_{N+1}$ и
        \[\idotsint\limits_K |f|\ dx_1 \dots dx_k\leq \idotsint\limits_{K_m\setminus K_{N+1}}|f|\ dx_1 \dots dx_k=a_m-a_{N+1}<\epsilon\]
        \item[$(\Leftarrow)$:]
        Выберем $\epsilon=1$ и положим $K_1=K_{\epsilon}$. Достроим $K_1$ до исчерпания $\{K_n\}$ области $\Omega$. По условию, если $K\subset \Omega\setminus K_{1}$, то
        \[\idotsint\limits_{K}|f|\ dx_1 \dots dx_k<1\]
        Заметим, что $K_n\setminus K_1\subset \Omega\setminus K_1$. Тогда для $n\geq 2$:
        \[\idotsint\limits_{K_n\setminus K_1}|f|\ dx_1 \dots dx_k<1\]
        Поскольку $K_1$ --- компакт и $f\in \mathcal{R_{\text{loc}}}(\Omega)$, то 
        \[\idotsint\limits_{K_1}|f|\ dx_1 \dots dx_k=C\]
        Таким образом
        \begin{multline*}
            \idotsint\limits_{K_n}|f|\ dx_1 \dots dx_k=\\
            =\idotsint\limits_{K_1}|f|\ dx_1 \dots dx_k+\idotsint\limits_{K_n\setminus K_1}|f|\ dx_1 \dots dx_k<C+1
        \end{multline*}
        Получили, что возрастающая последовательность интегралов
        \[\idotsint\limits_{K_n}|f|\ dx_1 \dots dx_k\]
        ограничена, а значит она сходится. Значит $|f|\in \widetilde{\mathcal{R}}(\Omega)$. Таким образом, по теореме \ref{kritintnes} выполнено $f\in \widetilde{\mathcal{R}}(\Omega)$.
    \end{itemize}
\end{proof}
\begin{example} [Интегрируемость $1/r^p$ по проколотому единичному шару и его внешности] $\\$
    Рассмотрим функцию 
    \[f(x_1,\dots,x_k)=\frac{1}{\left(\sqrt{x_1^2+\dots+x_k^2}\ \right)^p}\] 
    и интеграл
    \[\idotsint\limits_{0<\sqrt{x_1^2+\dots+x_k}<1}f(x_1,\dots,x_k)\ dx_1 \dots dx_k\]
    Поймем при каких $p$ он сходится: %Возьмем 
    %\[K_n=\frac{1}{n}\leq \sqrt{x_1^2+\dots+x_k^2}\leq 1-\frac{1}{n}\]
    %и перейдем к полярным координатам:
    \begin{multline*}
        \idotsint\limits_{0<\sqrt{x_1^2+\dots+x_k}<1}f(x_1,\dots,x_k)\ dx_1 \dots dx_k=\\
        =\idotsint\limits_{0<r<1}f(\phi(x_1,\dots,x_k))\cdot r^{k-1}\cdot \cos{\alpha_2}\cdot \dots \cos^{k-2}{\alpha_{k-1}}\ dr d\alpha_1 \dots d\alpha_{k-1}=\\
        =\int\limits_{0}^{1}\frac{1}{r^{p-k+1}}\ dr\cdot \int\limits_{0}^{2\pi}1 \ d\alpha_1\cdot \int\limits_{-\frac{\pi}{2}}^{\frac{\pi}{2}}\cos{\alpha_{2}}\ d\alpha_2 \dots \int\limits_{-\frac{\pi}{2}}^{\frac{\pi}{2}}\cos^{k-2}{\alpha_{k-1}}\ d\alpha_{k-1}\overset{p<k}{=}\\
        \overset{p<k}{=}\frac{1}{k-p}\cdot2\pi\cdot B\left(\frac{1+1}{2},\frac{0+1}{2}\right)\cdot B\left(\frac{2+1}{2},\frac{0+1}{2}\right)\dots B\left(\frac{k-1}{2},\frac{0+1}{2}\right)=\\
        =\frac{1}{k-p}\cdot2\pi\cdot \frac{\Gamma(1)\Gamma(\frac{1}{2})}{\Gamma(\frac{3}{2})}\cdot \frac{\Gamma(\frac{3}{2})\Gamma(\frac{1}{2})}{\Gamma(2)} \dots \frac{\Gamma(\frac{k-1}{2})\Gamma(\frac{1}{2})}{\Gamma(\frac{k}{2})}=\\
        =\frac{2\pi}{k-p}\cdot \frac{(\sqrt{\pi})^{k-2}}{\Gamma(\frac{k}{2})}=\frac{2\cdot \pi^{\frac{k}{2}}}{(k-p)\cdot \Gamma(\frac{k}{2})}
    \end{multline*}
    Теперь рассмотрим интеграл
    \[\idotsint\limits_{\sqrt{x_1^2+\dots+x_k}>1}f(x_1,\dots,x_k)\ dx_1 \dots dx_k\]
    \begin{multline*}
        \idotsint\limits_{K_n}f(x_1,\dots,x_k)\ dx_1 \dots dx_k=\\
        =\int\limits_{1}^{\infty}\frac{1}{r^{p-k+1}}\ dr\cdot \int\limits_{0}^{2\pi}1 \ d\alpha_1\cdot \int\limits_{-\frac{\pi}{2}}^{\frac{\pi}{2}}\cos{\alpha_{2}}\ d\alpha_2 \dots \int\limits_{-\frac{\pi}{2}}^{\frac{\pi}{2}}\cos^{k-2}{\alpha_{k-1}}\ d\alpha_{k-1}\overset{p>k}{=}\\
        \overset{p>k}{=}\frac{1}{p-k}\cdot2\pi\cdot B\left(\frac{1+1}{2},\frac{0+1}{2}\right)\cdot B\left(\frac{2+1}{2},\frac{0+1}{2}\right)\dots B\left(\frac{k-1}{2},\frac{0+1}{2}\right)=\\
        =\frac{1}{p-k}\cdot2\pi\cdot \frac{\Gamma(1)\Gamma(\frac{1}{2})}{\Gamma(\frac{3}{2})}\cdot \frac{\Gamma(\frac{3}{2})\Gamma(\frac{1}{2})}{\Gamma(2)} \dots \frac{\Gamma(\frac{k-1}{2})\Gamma(\frac{1}{2})}{\Gamma(\frac{k}{2})}=\\
        =\frac{2\pi}{p-k}\cdot \frac{(\sqrt{\pi})^{k-2}}{\Gamma(\frac{k}{2})}=\frac{2\cdot \pi^{\frac{k}{2}}}{(p-k)\cdot \Gamma(\frac{k}{2})}
    \end{multline*}
\end{example}
%\noindent \textit{Свойства аддитивности и интегрируемости по подмножеству в несобственных интегралов убираем из прораммы}
\newpage